\scanpage{44}
27.8.65

Mon Cher Serre,

Je continue à retourner dans tous les sens les questions de cycles algébriques; techniquement parlant je vois la question plus clairement, mais cependant il manque toujours le pas décisif. Comme je te connais allergique à la cohomologie, j'ai envie de t'énoncer les deux conjectures-clef, qu'il s'agit de démontrer, de façon purement géométrique i.e. sans parler de cohomologie. $X$ désigne une variété projective lisse connexe sur le corps algébriquement clos $k$, de dimension $n$, $Y$ une section hyperplane lisse. Par $C^i(X)$ je désigne le groupe des classes de cycles de codimension $i$ mod équivalence algébrique, tensorisé par $\underline{Q}$, par $\xi\in C^1(X)$ la classe de $Y$

\uline{Conjecture} A Pour tout entier $i$ tel que $2i\leq n$, le produit par $\xi^{n-2i}$ induit un isomorphisme $C^i(X)\xrightarrow{\sim}C^{n-i}(X)$.

Modulo vérification que je n'ai pas écrite en détail, il suffit même dans cette conjecture de vérifier la surjectivité, et de le faire pour $n=2i+1$. Le premier cas qui échappe pour le moment est $i=1,n=3$. Cette conjecture implique l'analogue, pour les $C^i$, des théorèmes connus de Lefschetz sur la comparaison de la cohomologie de $X$ et de $Y$ par les homomorphismes d'image directe et d'image inverse (pour la formulation correcte, se rappeler seulement que $C^i$ devient $H^{2i}$), en même temps qu'il implique les théorèmes de Lefschetz cohomologiques eux-mêmes, dont la forme formellement la plus forte consiste à affirmer que $\xi^{n-i}$ induit un isomorphisme de $H^i$ et $H^{2n-i}$ pour la cohomologie à coefficients $\underline{Q}_\ell$; je te signale que ce théorème purement cohomologique n'est pas même démontré pour le moment ! N'est démontré que le théorème de Lefschetz ``faible'' affirmant que la dimension cohomologique d'une variété affine de dimension $n$ est majorée par $n$, ce qui pour la variété affine $U=X-Y$ prend la forme : $H^i(X)\longrightarrow H^i(Y)$ est bijectif pour $i\leq n-2$, injectif pour $i=n-1$, ou encore : $H^i(Y)\longrightarrow H^{i+2}(X)$ est bijectif pour

\hfill $\ldots/\ldots$

\mspage{43}
